\documentclass[aspectratio=169,xcolor=dvipsnames]{beamer}

\AtBeginSection[]
{
    \begin{frame}
        \frametitle{Visão Geral}
        \tableofcontents[currentsection]
    \end{frame}
}

\usefonttheme{professionalfonts} % using non standard fonts for beamer
\usefonttheme{serif} % default family is serif
\usepackage{fontspec}
\setmainfont{XCharter}[
    UprightFont    = *-Roman,
    BoldFont       = *-Bold,
    ItalicFont     = *-Italic,
    BoldItalicFont = *-BoldItalic,
]

\usepackage{hyperref}
\usepackage{graphicx} % Allows including images
\usepackage{booktabs} % Allows the use of \toprule, \midrule and \bottomrule in tables
\input{./slides/SimplePlusTheme/beamerthemeSimplePlus.sty}
\input{./slides/SimplePlusTheme/beamerinnerthemeSimplePlus.sty}
\input{./slides/SimplePlusTheme/beamerfontthemeSimplePlus.sty}
\input{./slides/SimplePlusTheme/beamercolorthemeSimplePlus.sty}

\usepackage{listings}
\setbeamercovered{transparent}

\usepackage{biblatex}
\addbibresource{slides/reference.bib}
\renewcommand*{\bibfont}{\footnotesize}
% Tamanho dos slides de "Leituras Recomendadas": maior que a bibliografia
% completa (\bibfont acima, que rende \tiny via a redefinição de \footnotesize)
% para destacar as leituras curadas.
\newcommand{\leiturasize}{\small}

\usepackage[brazilian ]{babel}
\usepackage{csquotes}
\usepackage{pgfplots}
\pgfplotsset{compat=1.18}

\usepackage{emoji}
\usepackage{amsmath}
\usepackage[xcharter,bigdelims,vvarbb]{newtxmath}
% newtxmath's operators font requests the fontspec family `XCharter(0)' in OT1;
% without these substitutions math digits and operator names silently fall back
% to Computer Modern (LaTeX Font Warning: Font shape `OT1/XCharter(0)/m/n' undefined).
\DeclareFontFamily{OT1}{XCharter(0)}{}
\DeclareFontShape{OT1}{XCharter(0)}{m}{n}{<->ssub * XCharter-TLF/m/n}{}
\DeclareFontShape{OT1}{XCharter(0)}{m}{it}{<->ssub * XCharter-TLF/m/it}{}
\DeclareFontShape{OT1}{XCharter(0)}{b}{n}{<->ssub * XCharter-TLF/b/n}{}
\DeclareFontShape{OT1}{XCharter(0)}{bx}{n}{<->ssub * XCharter-TLF/b/n}{}

\usepackage{bm}
\usepackage[table]{xcolor}
\usepackage{xcolor}
\usepackage{algpseudocode}

\usepackage{cancel}
\usepackage{ulem}

\usetikzlibrary{shapes}
\usetikzlibrary{arrows}
\usetikzlibrary {arrows.meta}
\usetikzlibrary{calc,positioning}
\usepgfplotslibrary{fillbetween}
\usetikzlibrary{angles,quotes}
\usetikzlibrary{tikzmark}


\definecolor{PastelRed}{HTML}{FFC1CC}
\definecolor{PastelBlue}{HTML}{C2F0FF}
\definecolor{PastelBlue2}{HTML}{3182BD}

\definecolor{PastelPink}{HTML}{FFCBE1}
\definecolor{PastelGreen}{HTML}{D6E5BD}
\definecolor{PastelYellow}{HTML}{F9E1A8}
\definecolor{PastelBlue}{HTML}{BCD8EC}
\definecolor{PastelPurple}{HTML}{DCCCEC}
\definecolor{PastelOrange}{HTML}{FFDAB4}


\renewcommand{\footnotesize}{\tiny}

\newcommand{\mespor}{
  \ifcase\month
    \or janeiro
    \or fevereiro
    \or março
    \or abril
    \or maio
    \or junho
    \or julho
    \or agosto
    \or setembro
    \or outubro
    \or novembro
    \or dezembro
  \fi
}
% \usepackage{csquotes}
% \usepackage{minted}
% \usemintedstyle{xcode}
\definecolor{vscLightBackground}{HTML}{FFFFFF}
\definecolor{vscLightText}{HTML}{000000}
\definecolor{vscLightGreen}{HTML}{008000}        % For Comments AND Numbers
\definecolor{vscLightRed}{HTML}{A31515}          % For Strings
\definecolor{vscLightTeal}{HTML}{267F99}
\definecolor{vscBlue}{HTML}{0000FF}% For main keywords (if, for, def...)
\definecolor{vscLightPytorch}{HTML}{800080}      % For PyTorch keywords (dark magenta)

\lstdefinestyle{vscode-light}{
    language=Python,
    backgroundcolor=\color{vscLightBackground},
    basicstyle=\ttfamily\scriptsize\color{vscLightText},
    keywordstyle=\color{vscLightTeal},
    commentstyle=\color{vscLightGreen},
    stringstyle=\color{vscLightRed},
    % Style for PyTorch keywords (using keyword group 2)
    keywordstyle=[2]\color{vscLightPytorch}, 
    morekeywords={self, True, False, None},
    morekeywords=[2]{
        torch, Tensor, device, dtype, to, nn, optim, autograd, utils, data,
        Module, Linear, Conv2d, ReLU, CrossEntropyLoss, Adam, SGD,
        Dataset, DataLoader, randn, zeros, ones, arange, cat, stack,
        view, reshape, sum, mean, max, min, matmul, no_grad, backward
    },
    keywordstyle=[3]\color{vscBlue}, 
    morekeywords=[3]{
        model, loss, y, y_hat, x , optimizer
    },
    frame=single,
    tabsize=1,
    showstringspaces=false,
    breaklines=true,
    captionpos=b,
    extendedchars=true,
    numbers=left,
}

\usepackage[cache=true]{minted}
% minted v3+ (TeX Live 2024+) auto-detects the output directory.
% minted v2 (Ubuntu apt TeX Live 2023) needs it set explicitly to
% match latexmkrc's $out_dir = 'build'.
\makeatletter
\@ifpackagelater{minted}{2024/01/01}{}{%
  \def\minted@outputdir{build/}%
}
\makeatother
\setminted{
    fontsize=\scriptsize,
    style=vs,
    breaklines=true,
    numbers=left,
    numbersep=5pt,
    frame=single,
}
\providecommand{\citeauthoryear}[1]{(\citeauthor{#1},~\citeyear{#1})}

\providecommand{\thetav}{\bm{\theta}}

% vetor coluna 2D com colchetes
\providecommand{\cvec}[2]{\begin{bmatrix} #1 \\ #2 \end{bmatrix}}

\providecommand{\varvector}[1]{\mathbf{#1}}
% instance
\providecommand{\X}{\varvector{X}}
\providecommand{\mlinstance}{\varvector{X}}
\providecommand{\mlinstancei}{\varvector{X}^{(i)}}
\providecommand{\mlinstancex}[1]{\varvector{X}^{(#1)}}

% target
\providecommand{\y}{\varvector{y}}
\providecommand{\mltarget}{\varvector{y}}
\providecommand{\mltargeti}{\varvector{y}^{(i)}}
\providecommand{\mltargetx}[1]{\varvector{y}^{(#1)}}

%hats
\providecommand{\yhat}{\hat{\varvector{y}}}
\providecommand{\mltargethat}{\hat{\varvector{y}}}
\providecommand{\mltargethati}{\hat{\varvector{y}}^{(i)}}
\providecommand{\mltargethatx}[1]{\hat{\varvector{y}}^{(#1)}}

% linear reg
\providecommand{\mllinearregression}{\mltargethat=\mlinstance\bm{\theta}}

\providecommand{\tikzarrowcolor}[3]{%
  \begin{tikzpicture}[remember picture, overlay]
    % Arrow body
    \path[fill=#3] ([shift={(#1,#2)}]current page.south west) 
      rectangle ++(0.3, -0.2);
    % Arrow head
    \path[fill=#3] ([shift={(#1+0.3,#2+0.1)}]current page.south west) -- 
         ++(0, -0.4) -- 
         ++(0.2, 0.2) -- 
         cycle;
  \end{tikzpicture}%
}

\providecommand{\tikzarrow}[2]{%
    \tikzarrowcolor{#1}{#2}{Red}
}



\DeclareMathOperator*{\argmax}{argmax}
\DeclareMathOperator*{\argmin}{argmin}



%----------------------------------------------------------------------------------------
%    TITLE PAGE
%----------------------------------------------------------------------------------------

\title{Deep Learning}
\subtitle{Overfitting, Underfitting e Regularização}

\author{Lucas Silveira Kupssinskü}

\institute
{
    Machine Learning Theory and Applications Laboratory \\
    PUCRS
}
\date{\mespor de \the\year}

%----------------------------------------------------------------------------------------
%    PRESENTATION SLIDES
%----------------------------------------------------------------------------------------

\begin{document}

  \begin{frame}
    \titlepage
  \end{frame}

  \begin{frame}{Visão Geral}
    \tableofcontents
  \end{frame}

  %======================================================================
  \section{Motivação: overfitting na prática}
  %======================================================================

  %--- Slide: MNIST 1D ---
  \begin{frame}{Conjunto de dados: MNIST 1D\footfullcite{greydanus2020scaling}}
    \begin{columns}
      \begin{column}{.55\linewidth}
        \begin{itemize}
          \item Versão simplificada do MNIST, projetada para experimentos
            rápidos com redes neurais.
          \item Cada classe possui um \textit{template} unidimensional
            que sofre transformações estocásticas.
          \item Pipeline de geração:
            \begin{enumerate}
              \item \textit{Templates} para as 10 classes.
              \item Distorções aleatórias (translação, escala, espelhamento).
              \item Adição de ruído.
              \item Subamostragem em 40 pontos.
            \end{enumerate}
          \item Repositório: \href{https://github.com/greydanus/mnist1d}{\textit{github.com/greydanus/mnist1d}}.
        \end{itemize}
      \end{column}
      \begin{column}{.45\linewidth}
        \centering
        \includegraphics[width=\linewidth,height=.7\textheight,keepaspectratio]{slides/img/01-mlp/reg-mnist1d-templates.png}\\
        \footnotesize a) \textit{Templates}, b) exemplos de treino, c) ruído, d) subamostragem em 40 pontos.
      \end{column}
    \end{columns}
  \end{frame}

  %--- Slide: setup do MLP ---
  \begin{frame}{MLP usado no experimento}
    \begin{columns}
      \begin{column}{.55\linewidth}
        \begin{itemize}
          \item Entradas: 40 (uma por ponto amostrado).
          \item Saída: 10 logits, ativação \textit{softmax}.
          \item Duas camadas ocultas, 100 unidades cada, ativação ReLU.
          \item Conjunto de treino: 4000 instâncias.
          \item Otimizador: SGD, \textit{batch size} 100, \textit{learning rate} 0.1.
          \item Total de iterações: 6000, equivalente a:
            $$\text{épocas} = \frac{6000 \cdot 100}{4000} = 150.$$
        \end{itemize}
      \end{column}
      \begin{column}{.45\linewidth}
        \centering
        \begin{tikzpicture}[
          node distance=0.55cm and 0.9cm,
          every node/.style={font=\tiny},
          unit/.style={circle,draw,minimum size=0.32cm,inner sep=0pt}
        ]
          \node[unit] (i1) {};
          \node[unit, below=of i1] (i2) {};
          \node[unit, below=of i2] (i3) {};
          \node[below=0.05cm of i3,font=\tiny] {$x_{40}$};
          \node[above=0.05cm of i1,font=\tiny] {$x_1$};

          \node[unit, right=of i1] (h11) {};
          \node[unit, below=of h11] (h12) {};
          \node[unit, below=of h12] (h13) {};

          \node[unit, right=of h11] (h21) {};
          \node[unit, below=of h21] (h22) {};
          \node[unit, below=of h22] (h23) {};

          \node[unit, right=of h22] (o) {};
          \node[right=0.05cm of o,font=\tiny] {$\hat{\mathbf{y}}$};

          \foreach \i in {i1,i2,i3}
            \foreach \j in {h11,h12,h13} \draw[->,gray!60] (\i) -- (\j);
          \foreach \i in {h11,h12,h13}
            \foreach \j in {h21,h22,h23} \draw[->,gray!60] (\i) -- (\j);
          \foreach \i in {h21,h22,h23} \draw[->,gray!60] (\i) -- (o);
        \end{tikzpicture}\\
        \footnotesize Esquema (representação reduzida).
      \end{column}
    \end{columns}
  \end{frame}

  %--- Slide: resultados de overfitting ---
  \begin{frame}{Resultados do treinamento}
    \begin{columns}
      \begin{column}{.55\linewidth}
        \begin{itemize}
          \item A taxa de erro no \textit{conjunto de treino} cai a zero.
          \item A função de custo no \textit{conjunto de teste} primeiro diminui,
            mas a partir de certo ponto começa a aumentar.
          \item O modelo continua se ajustando ao ruído do treino e
            \alert{generaliza pior}, comportamento clássico de \textit{overfitting}.
          \item Capacidade alta sem regularização leva a esse padrão.
        \end{itemize}
      \end{column}
      \begin{column}{.45\linewidth}
        \centering
        \includegraphics[width=\linewidth,height=.65\textheight,keepaspectratio]{slides/img/01-mlp/reg-mnist1d-overfit.png}
      \end{column}
    \end{columns}
  \end{frame}

  %--- Slide: por que ocorre overfitting ---
  \begin{frame}{Por que isso ocorre? Toy MLP de 3 unidades}
    \begin{columns}
      \begin{column}{.42\linewidth}
        \begin{itemize}
          \item MLP com uma camada oculta de 3 unidades ReLU.
          \item Parâmetros nas linhas pontilhadas estão \alert{congelados}.
          \item Variando $\bm{\phi} = \{\beta, \omega_1, \omega_2, \omega_3\}$
            obtemos diferentes funções lineares por partes.
          \item Em e), f) e g): parametrizações distintas geram saídas distintas,
            mostrando como diferentes mínimos do treino correspondem a
            funções muito diferentes.
        \end{itemize}
      \end{column}
      \begin{column}{.58\linewidth}
        \centering
        \includegraphics[width=\linewidth,height=.78\textheight,keepaspectratio]{slides/img/01-mlp/reg-toy-variance.png}
      \end{column}
    \end{columns}
  \end{frame}

  %======================================================================
  \section{Fontes de erro: ruído, viés e variância}
  %======================================================================

  %--- Slide: visão geral das três fontes ---
  \begin{frame}{Três fontes de erro de teste}
    \centering
    \includegraphics[width=.95\linewidth,height=.78\textheight,keepaspectratio]{slides/img/01-mlp/reg-sources-overview.png}
  \end{frame}

  %--- Slide: ruído ---
  \begin{frame}{Ruído}
    \begin{columns}
      \begin{column}{.5\linewidth}
        \begin{itemize}
          \item Mesmo recuperando exatamente a função geradora dos dados
            (em preto), o erro não vai a zero.
          \item Origem: fatores estocásticos no processo gerador, ou
            variáveis observadas que não capturam toda a informação relevante.
          \item Componente \alert{irredutível}: nenhum modelo elimina.
        \end{itemize}
      \end{column}
      \begin{column}{.5\linewidth}
        \centering
        \includegraphics[width=\linewidth,height=.7\textheight,keepaspectratio]{slides/img/01-mlp/reg-source-noise.png}
      \end{column}
    \end{columns}
  \end{frame}

  %--- Slide: viés ---
  \begin{frame}{Viés}
    \begin{columns}
      \begin{column}{.5\linewidth}
        \begin{itemize}
          \item Mesmo com os melhores parâmetros possíveis, o MLP em ciano
            não modela perfeitamente a função em preto.
          \item Diferença entre a função que o modelo é \textit{capaz} de
            representar e a função real.
          \item Todo modelo carrega um \textit{viés indutivo}: sem ele, não há
            generalização possível.
        \end{itemize}
      \end{column}
      \begin{column}{.5\linewidth}
        \centering
        \includegraphics[width=\linewidth,height=.7\textheight,keepaspectratio]{slides/img/01-mlp/reg-source-bias.png}
      \end{column}
    \end{columns}
  \end{frame}

  %--- Slide: variância ---
  \begin{frame}{Variância}
    \begin{columns}
      \begin{column}{.5\linewidth}
        \begin{itemize}
          \item Ao treinar com um número limitado de instâncias, não recuperamos
            necessariamente o melhor ajuste à função geradora.
          \item Pequenas mudanças no conjunto de treino produzem modelos
            diferentes.
          \item Variância mede o quanto o modelo aprendido oscila com o
            \textit{dataset} de treino.
        \end{itemize}
      \end{column}
      \begin{column}{.5\linewidth}
        \centering
        \includegraphics[width=\linewidth,height=.7\textheight,keepaspectratio]{slides/img/01-mlp/reg-source-variance.png}
      \end{column}
    \end{columns}
  \end{frame}

  %--- Slide: derivação bias-variance, setup ---
  \begin{frame}{Decomposição viés-variância: \textit{setup}}
    \small
    Considere regressão escalar, com:
    \begin{itemize}
      \item Função geradora $\mu(x) = \mathbb{E}[y \mid x]$.
      \item Observações $y = \mu(x) + \epsilon$, em que $\mathbb{E}[\epsilon] = 0$ e $\text{Var}(\epsilon) = \sigma^2$.
      \item Predição $f(x;\mathcal{D})$ treinada em um \textit{dataset} $\mathcal{D}$.
      \item Custo quadrático $\ell = (y - f(x;\mathcal{D}))^2$.
    \end{itemize}

    \smallskip
    Queremos decompor o erro esperado em $x$:
    $$\mathbb{E}_{\mathcal{D},\,y}\!\left[(y - f(x;\mathcal{D}))^2\right].$$

    \smallskip
    Estratégia: dois passos de soma e subtração de um termo conveniente.
    Primeiro $\mu(x)$, isolando o ruído. Depois o modelo médio
    $\bar f(x) = \mathbb{E}_{\mathcal{D}}[f(x;\mathcal{D})]$, separando viés e variância.
  \end{frame}

  %--- Slide: derivação bias-variance, passo 1 ---
  \begin{frame}{Passo 1: isolar o ruído}
    \small
    Fixamos $f(x)$ e tomamos a esperança apenas em $y$.
    Somando e subtraindo $\mu(x)$:
    $$y - f(x) = \bigl(y - \mu(x)\bigr) + \bigl(\mu(x) - f(x)\bigr).$$

    \onslide<2->{%
      Elevando ao quadrado:
      $$(y - f(x))^2 = (y - \mu(x))^2 \,+\, 2\,(y - \mu(x))(\mu(x) - f(x)) \,+\, (\mu(x) - f(x))^2.$$
    }

    \onslide<3->{%
      Aplicando $\mathbb{E}_y$ e usando que $(\mu(x) - f(x))$ é constante em $y$, o termo cruzado se anula:
      $$\mathbb{E}_y\!\left[(y - \mu(x))(\mu(x) - f(x))\right]
        = (\mu(x) - f(x))\,\underbrace{\mathbb{E}_y[y - \mu(x)]}_{=\,0} = 0.$$
    }

    \onslide<4->{%
      Restam dois termos:
      $$\mathbb{E}_y[(y - f(x))^2]
        = \underbrace{\mathbb{E}_y[(y - \mu(x))^2]}_{\sigma^2,\ \text{ruído}}
        \,+\, \underbrace{(\mu(x) - f(x))^2}_{\text{erro do modelo}}.$$
    }
  \end{frame}

  %--- Slide: derivação bias-variance, passo 2 ---
  \begin{frame}{Passo 2: separar viés e variância}
    \small
    O termo de erro do modelo ainda depende do \textit{dataset} via $f(x;\mathcal{D})$.
    Tomamos agora a esperança sobre $\mathcal{D}$, somando e subtraindo $\bar f(x)$:
    $$\mu(x) - f(x;\mathcal{D}) = \bigl(\mu(x) - \bar f(x)\bigr) + \bigl(\bar f(x) - f(x;\mathcal{D})\bigr).$$

    \onslide<2->{%
      Elevando ao quadrado:
      $$(\mu - f)^2 = (\mu - \bar f)^2 \,+\, 2\,(\mu - \bar f)(\bar f - f) \,+\, (\bar f - f)^2.$$
    }

    \onslide<3->{%
      Aplicando $\mathbb{E}_{\mathcal{D}}$, o termo cruzado se anula porque $(\mu - \bar f)$ é constante em $\mathcal{D}$ e $\mathbb{E}_{\mathcal{D}}[\bar f - f] = \bar f - \bar f = 0$:
      $$\mathbb{E}_{\mathcal{D}}\!\left[(\mu - \bar f)(\bar f - f)\right]
        = (\mu - \bar f)\,\underbrace{\mathbb{E}_{\mathcal{D}}[\bar f - f]}_{=\,0} = 0.$$
    }

    \onslide<4->{%
      Restam dois termos:
      $$\mathbb{E}_{\mathcal{D}}[(\mu(x) - f(x;\mathcal{D}))^2]
        = \underbrace{(\mu(x) - \bar f(x))^2}_{\text{viés}^2}
        \,+\, \underbrace{\mathbb{E}_{\mathcal{D}}[(f(x;\mathcal{D}) - \bar f(x))^2]}_{\text{variância}}.$$
    }
  \end{frame}

  %--- Slide: derivação bias-variance, resultado ---
  \begin{frame}{Resultado: três fontes de erro}
    \small
    Combinando os dois passos, o erro esperado em $x$ se decompõe em três termos:
    $$\mathbb{E}_{\mathcal{D},\,y}\!\left[(y - f(x;\mathcal{D}))^2\right]
      = \underbrace{\sigma^2}_{\text{ruído}}
      \,+\, \underbrace{(\mu(x) - \bar f(x))^2}_{\text{viés}^2}
      \,+\, \underbrace{\mathbb{E}_{\mathcal{D}}[(f(x;\mathcal{D}) - \bar f(x))^2]}_{\text{variância}}.$$

    \begin{itemize}
      \item \textbf{Ruído} $\sigma^2$: componente irredutível, depende apenas do processo gerador.
      \item \textbf{Viés}$^2$: distância entre o modelo médio $\bar f$ e a função real $\mu$.
      \item \textbf{Variância}: oscilação do modelo aprendido em torno de $\bar f$, induzida pela amostra de treino.
    \end{itemize}

    \onslide<2->{%
      Esta decomposição aditiva vale exatamente para o \alert{custo quadrático},
      qualquer que seja o modelo $f$. Para outras funções de custo (0-1, entropia
      cruzada) ela não é tão limpa, mas a intuição de viés e variância permanece útil.
    }
  \end{frame}

  %======================================================================
  \section{Trade-off viés-variância e double descent}
  %======================================================================

  %--- Slide: dataset size vs variância ---
  \begin{frame}{Tamanho do \textit{dataset} vs. variância}
    \begin{columns}
      \begin{column}{.45\linewidth}
        \begin{itemize}
          \item Mesmo modelo (capacidade fixa: 3 regiões lineares).
          \item Com 6 instâncias: ajustes muito variáveis entre \textit{datasets}.
          \item Com 100 instâncias: ajustes próximos uns dos outros.
          \item Mais dados, menor variância.
        \end{itemize}
      \end{column}
      \begin{column}{.55\linewidth}
        \centering
        \includegraphics[width=\linewidth,height=.78\textheight,keepaspectratio]{slides/img/01-mlp/reg-dataset-vs-variance.png}
      \end{column}
    \end{columns}
  \end{frame}

  %--- Slide: capacidade do modelo vs variância ---
  \begin{frame}{Capacidade do modelo vs. variância}
    \begin{columns}
      \begin{column}{.45\linewidth}
        \begin{itemize}
          \item Mesmo \textit{dataset}, modelos com mais regiões lineares.
          \item Mais capacidade, ajustes mais sensíveis à amostra de treino.
          \item Mais capacidade, maior variância.
        \end{itemize}
      \end{column}
      \begin{column}{.55\linewidth}
        \centering
        \includegraphics[width=\linewidth,height=.78\textheight,keepaspectratio]{slides/img/01-mlp/reg-capacity-vs-variance.png}
      \end{column}
    \end{columns}
  \end{frame}

  %--- Slide: trade-off bias-variance clássico ---
  \begin{frame}{Trade-off viés-variância clássico}
    \begin{columns}
      \begin{column}{.45\linewidth}
        \begin{itemize}
          \item Pouca capacidade: viés alto domina o erro.
          \item Muita capacidade: variância alta domina o erro.
          \item Existe um ponto intermediário que minimiza o erro total.
          \item Esta é a visão \textit{clássica}, antes das redes neurais
            superparametrizadas.
        \end{itemize}
      \end{column}
      \begin{column}{.55\linewidth}
        \centering
        \includegraphics[width=\linewidth,height=.78\textheight,keepaspectratio]{slides/img/01-mlp/reg-bv-tradeoff.png}
      \end{column}
    \end{columns}
  \end{frame}

  %--- Slide: double descent (overlay com 3 datasets) ---
  \begin{frame}{Double descent em redes neurais\footfullcite{belkin2019reconciling}}
    \begin{columns}
      \begin{column}{.4\linewidth}
        \begin{itemize}
          \item Em redes neurais o erro de teste pode \alert{voltar a cair}
            depois de passar pelo limiar de interpolação.
          \item Esse limiar ocorre quando o número de parâmetros se aproxima
            do número de instâncias de treino.
          \item Padrão observado em vários conjuntos:
            \only<1>{MNIST-1D.}\only<2>{MNIST-1D, MNIST.}\only<3->{MNIST-1D, MNIST, CIFAR-100.}
        \end{itemize}
      \end{column}
      \begin{column}{.6\linewidth}
        \centering
        \only<1>{\includegraphics[width=\linewidth,height=.75\textheight,keepaspectratio]{slides/img/01-mlp/reg-dd-mnist1d.png}}
        \only<2>{\includegraphics[width=\linewidth,height=.75\textheight,keepaspectratio]{slides/img/01-mlp/reg-dd-mnist.png}}
        \only<3->{\includegraphics[width=\linewidth,height=.75\textheight,keepaspectratio]{slides/img/01-mlp/reg-dd-cifar100.png}}
      \end{column}
    \end{columns}
  \end{frame}

  %--- Slide: ilustração double descent (Belkin) ---
  \begin{frame}{Por que ocorre? Hipóteses}
    \begin{columns}
      \begin{column}{.45\linewidth}
        \begin{itemize}
          \item Talvez a inicialização favoreça funções mais suaves.
          \item Talvez o treinamento favoreça funções mais suaves.
          \item Ainda é um tema aberto, com bons palpites mas sem consenso.
          \item No regime moderno, a relação capacidade $\times$ erro de teste
            tem dois mínimos: o ``sweet spot'' clássico e o regime
            superparametrizado.
        \end{itemize}
      \end{column}
      \begin{column}{.55\linewidth}
        \centering
        \includegraphics[width=\linewidth,height=.7\textheight,keepaspectratio]{slides/img/01-mlp/reg-belkin.jpeg}\\
        \footnotesize \citeauthoryear{belkin2019reconciling}
      \end{column}
    \end{columns}
  \end{frame}

  %======================================================================
  \section{Regularização explícita e empírica}
  %======================================================================

  %--- Slide: taxonomia ---
  \begin{frame}{Tipos de regularização}
    \begin{itemize}
      \item \textbf{Explícita}: termos extras na função de custo,
        L1 e L2.
      \item \textbf{Empírica}: técnicas que modificam o procedimento de
        treinamento, \textit{early stopping}, \textit{dropout},
        \textit{data augmentation}, \textit{label smoothing}.
      \item \textbf{Implícita}: efeitos de regularização que surgem
        do próprio otimizador (descida de gradiente e SGD).
      \item Todas têm o mesmo objetivo: reduzir variância sem aumentar
        demais o viés.
    \end{itemize}
  \end{frame}

  %--- Slide: L1 e L2 (formulação) ---
  \begin{frame}{Regularização L1 e L2}
    \begin{itemize}
      \item Adicionamos um termo de penalidade à função de custo:
        $$\tilde{\mathcal{L}}(\bm{\theta}) = \mathcal{L}(\bm{\theta}) + \lambda \, \Omega(\bm{\theta}),
          \qquad \Omega_{L2} = \|\bm{\theta}\|_2^2,
          \qquad \Omega_{L1} = \|\bm{\theta}\|_1.$$
      \item Enviesamos a otimização em favor de parâmetros de \alert{baixa magnitude}.
      \item Visão probabilística: equivale a adicionar um \textit{prior}
        no critério de máxima verossimilhança.
      \begin{itemize}
        \item L2 corresponde a um \textit{prior} gaussiano,
          $\bm{\theta} \sim \mathcal{N}(0, \tau^2 \mathbf{I})$.
        \item L1 corresponde a um \textit{prior} laplaciano,
          que induz \alert{esparsidade}.
      \end{itemize}
      \item Como redes neurais atuam em problemas variados, o \textit{prior}
        precisa ser genérico.
    \end{itemize}
  \end{frame}

  %--- Slide: L2 como decaimento de pesos ---
  \begin{frame}{L2 e \textit{weight decay}}
    \small
    Para L2, a atualização SGD se reescreve como:
    \begin{align*}
      \bm{\theta}_{t+1}
        &= \bm{\theta}_t - \eta \, \nabla_{\bm{\theta}}\!\left[\mathcal{L}(\bm{\theta}_t) + \lambda \|\bm{\theta}_t\|_2^2\right]\\
        &= \bm{\theta}_t - \eta \, \nabla_{\bm{\theta}} \mathcal{L}(\bm{\theta}_t) - 2\eta\lambda\,\bm{\theta}_t\\
        &= \underbrace{(1 - 2\eta\lambda)\,\bm{\theta}_t}_{\text{decaimento}}
           \; - \; \eta \,\nabla_{\bm{\theta}} \mathcal{L}(\bm{\theta}_t).
    \end{align*}
    \begin{itemize}
      \item A cada passo, os pesos encolhem multiplicativamente por $(1 - 2\eta\lambda)$.
      \item Esse comportamento dá origem ao nome \textit{weight decay}.
    \end{itemize}
    \vspace{0.3em}
    {\footnotesize\color{gray!50!black}No PyTorch, o argumento \texttt{weight\_decay} corresponde ao coeficiente já em frente de $\bm\theta$ no \textit{update} (i.e., $\lambda_{\text{PT}} = 2\lambda$ desta derivação).}
  \end{frame}

  %--- Slide: efeito de aumentar lambda ---
  \begin{frame}{Efeito de aumentar $\lambda$}
    \begin{columns}
      \begin{column}{.4\linewidth}
        \begin{itemize}
          \item $\lambda$ pequeno: o ajuste se aproxima do caso sem regularização,
            modelo sensível ao ruído dos dados.
          \item $\lambda$ grande: o ajuste se torna mais suave,
            modelo se afasta dos pontos individuais.
          \item Existe um $\lambda$ intermediário que generaliza melhor.
        \end{itemize}
      \end{column}
      \begin{column}{.6\linewidth}
        \centering
        \includegraphics[width=\linewidth,height=.78\textheight,keepaspectratio]{slides/img/01-mlp/reg-lambda-effect.png}
      \end{column}
    \end{columns}
  \end{frame}

  %--- Slide: L1 vs L2 (paisagem de loss) ---
  \begin{frame}{L1 vs. L2 na paisagem de \textit{loss}}
    \begin{columns}
      \begin{column}{.4\linewidth}
        \begin{itemize}
          \item Painel a): \textit{loss} original com vários mínimos locais.
          \item Painel b): termo de regularização sozinho, com mínimo único na origem.
          \item Painel c): soma das duas paisagens, com poucos mínimos próximos da origem.
          \item L1 favorece componentes \alert{exatamente zero} (os vértices do losango
            do L1 ficam sobre os eixos); L2 favorece magnitudes pequenas mas não zero.
        \end{itemize}
      \end{column}
      \begin{column}{.6\linewidth}
        \centering
        \includegraphics[width=\linewidth,height=.78\textheight,keepaspectratio]{slides/img/01-mlp/reg-l1-vs-l2.png}
      \end{column}
    \end{columns}
  \end{frame}

  %--- Slide: early stopping (conceito) ---
  \begin{frame}{\textit{Early stopping}}
    \begin{itemize}
      \item Monitoramos o desempenho da rede em um conjunto de validação
        e interrompemos o treinamento antes da \textit{loss} de validação
        voltar a subir.
      \item Tem efeito \alert{análogo} a L2, pois impede que os pesos
        cresçam indiscriminadamente.
      \item Hiperparâmetro principal: \textit{patience}, número de iterações
        sem melhoria que toleramos antes de parar.
      \item Bônus: também economiza tempo de computação.
    \end{itemize}
  \end{frame}

  %--- Slide: early stopping (gráfico) ---
  \begin{frame}{\textit{Early stopping} em ação}
    \centering
    \includegraphics[width=.85\linewidth,height=.78\textheight,keepaspectratio]{slides/img/01-mlp/reg-early-stopping.png}
  \end{frame}

  %--- Slide: dropout (conceito + diagrama) ---
  \begin{frame}{\textit{Dropout}\footfullcite{srivastava2014dropout}}
    \begin{columns}
      \begin{column}{.45\linewidth}
        \begin{itemize}
          \item Durante o treinamento, unidades de uma camada oculta são
            \alert{desligadas} aleatoriamente.
          \item Equivale a multiplicar as ativações por uma máscara binária
            $\mathbf{m} \sim \text{Bernoulli}(1 - p)$.
          \item A cada \textit{forward}, uma sub-rede diferente é ativada,
            efeito similar a treinar um \textit{ensemble} implícito.
        \end{itemize}
      \end{column}
      \begin{column}{.55\linewidth}
        \centering
        \includegraphics[width=\linewidth,height=.6\textheight,keepaspectratio]{slides/img/01-mlp/reg-dropout-net.png}\\
        \footnotesize Esquerda: rede densa. Direita: rede com unidades desligadas.
      \end{column}
    \end{columns}
  \end{frame}

  %--- Slide: dropout (efeito visual progressivo) ---
  \begin{frame}{\textit{Dropout}: por que ajuda}
    \begin{columns}
      \begin{column}{.4\linewidth}
        \begin{itemize}
          \item a) Sem \textit{dropout} o modelo já interpola os dados de treino,
            mas com um pico isolado.
          \item b) Ao desligar uma unidade, esse pico vira um aumento brusco
            na \textit{loss}.
          \item c) Continuando a desligar unidades aleatoriamente, o pico
            é gradualmente eliminado.
        \end{itemize}
      \end{column}
      \begin{column}{.6\linewidth}
        \centering
        \includegraphics[width=\linewidth,height=.78\textheight,keepaspectratio]{slides/img/01-mlp/reg-dropout-effect.png}
      \end{column}
    \end{columns}
  \end{frame}

  %--- Slide: dropout em inferência ---
  \begin{frame}{\textit{Dropout} em inferência}
    \begin{itemize}
      \item Em treino, cada unidade é desligada com probabilidade $p$, então a magnitude
        esperada das ativações é $(1 - p)$ vezes a magnitude original.
      \item Em inferência precisamos compensar esse fator. Duas opções:
      \begin{itemize}
        \item \textbf{Rescaling determinístico}: multiplicar as ativações
          de teste por $(1 - p)$, ou equivalentemente dividi-las em treino
          por $(1 - p)$ (\textit{inverted dropout}, padrão em PyTorch e Keras).
        \item \textbf{Monte Carlo Dropout}\footfullcite{gal2016dropout}: rodar o \textit{forward}
          múltiplas vezes com máscaras diferentes e combinar as saídas,
          análogo a um \textit{ensemble} explícito; também fornece uma
          estimativa de incerteza.
      \end{itemize}
    \end{itemize}
  \end{frame}

  %--- Slide: data augmentation ---
  \begin{frame}{\textit{Data augmentation}}
    \begin{columns}
      \begin{column}{.5\linewidth}
        \begin{itemize}
          \item Geramos novas instâncias aplicando transformações nos dados
            originais que preservam o rótulo.
          \item Em imagens: rotações, \textit{flip}, mudança de cor,
            equalização, \textit{random crop}, ruído.
          \item Codificamos \alert{invariâncias} desejadas no conjunto de treino,
            sem alterar a arquitetura.
          \item Efeito: aumenta a diversidade efetiva do \textit{dataset},
            reduzindo a variância do modelo.
        \end{itemize}
      \end{column}
      \begin{column}{.5\linewidth}
        \centering
        \includegraphics[width=.45\linewidth,height=.45\textheight,keepaspectratio]{slides/img/01-mlp/reg-augmentation-1.png}\\
        \vspace{0.2em}
        \includegraphics[width=\linewidth,height=.32\textheight,keepaspectratio]{slides/img/01-mlp/reg-augmentation-2.png}\\
        \footnotesize Original (acima) e variantes geradas (abaixo).
      \end{column}
    \end{columns}
  \end{frame}

  %--- Slide: label smoothing (conceito) ---
  \begin{frame}{\textit{Label smoothing}\footfullcite{szegedy2016rethinking}}
    \begin{itemize}
      \item \textit{Data augmentation} transforma as \alert{entradas};
        podemos também regularizar transformando os \alert{alvos}.
      \item Com alvo \textit{one-hot}, a entropia cruzada só é minimizada
        quando a probabilidade da classe correta tende a 1, o que exige
        \textit{logits} cada vez maiores: o modelo é empurrado para o
        \alert{\textit{overconfidence}}.
      \item \textit{Label smoothing} redistribui uma fração $\alpha$ da massa
        de probabilidade uniformemente entre as $C$ classes:
        $$\mathbf{y}^{(i)}_{ls} = (1 - \alpha)\,\mathbf{y}^{(i)}_{oh} + \frac{\alpha}{C}.$$
      \item Casos limite: $\alpha = 0$ mantém o alvo \textit{one-hot} original;
        $\alpha = 1$ produz um alvo uniforme, sem informação de classe.
      \item Na prática, valores pequenos como $\alpha = 0{,}1$ são os mais comuns.
    \end{itemize}
  \end{frame}

  %--- Slide: label smoothing (efeito) ---
  \begin{frame}{\textit{Label smoothing}: por que ajuda\footfullcite{muller2019labelsmoothing}}
    \begin{itemize}
      \item Exemplo com $C = 4$, classe correta 3 e $\alpha = 0{,}1$:
        $$\mathbf{y}_{oh} = (0,\; 0,\; 1,\; 0)
          \quad\longrightarrow\quad
          \mathbf{y}_{ls} = (0{,}025,\; 0{,}025,\; 0{,}925,\; 0{,}025).$$
      \item O alvo suavizado é atingível com \textit{logits} \alert{finitos}:
        a diferença entre o \textit{logit} da classe correta e os demais
        para de crescer, contendo a magnitude dos pesos
        (efeito análogo ao de L2 na camada de saída).
      \item Modelos treinados com \textit{label smoothing} tendem a ficar
        \alert{melhor calibrados}: a confiança predita se aproxima da
        acurácia real.
      \item Custo: as representações da penúltima camada perdem informação
        sobre a semelhança entre classes, o que prejudica, por exemplo,
        a destilação de conhecimento para um modelo menor.
    \end{itemize}
  \end{frame}

  %======================================================================
  \section{Regularização implícita}
  %======================================================================

  %--- Slide: conceito ---
  \begin{frame}{Regularização implícita}
    \begin{itemize}
      \item Fato peculiar: a descida de gradiente encontra soluções que
        \alert{generalizam}, mesmo em redes muito superparametrizadas.
      \item Tanto o \textit{batch GD} quanto o SGD não fazem uma trajetória
        direta para o mínimo local da função de custo: existe uma preferência
        por algumas soluções em detrimento de outras.
      \item Essas preferências, que não são impostas explicitamente na
        função de custo, recebem o nome de \alert{regularização implícita}.
    \end{itemize}
  \end{frame}

  %--- Slide: implicit gradient regularization (Barrett & Dherin) ---
  \begin{frame}{\textit{Implicit Gradient Regularization}\footfullcite{barrett2020implicit}}
    \begin{itemize}
      \item Versão contínua da descida de gradiente:
        $\dot{\bm{\theta}}(t) = -\nabla \mathcal{L}(\bm{\theta}(t))$.
      \item A versão discreta (passos de tamanho $\eta$) aproxima a versão contínua
        de uma função de custo \alert{modificada}:
        $$\tilde{\mathcal{L}}_{GD}(\bm{\theta}) = \mathcal{L}(\bm{\theta})
          + \frac{\eta}{4} \left\| \nabla \mathcal{L}(\bm{\theta}) \right\|^2.$$
      \item Consequência: a solução é \alert{repelida} de regiões em que a norma
        do gradiente é grande (alta declividade da \textit{loss}).
      \item GD prefere bacias de mínimo \textit{rasas} e largas, que tendem
        a generalizar melhor.
    \end{itemize}
  \end{frame}

  %--- Slide: SGD favorece gradiente estável ---
  \begin{frame}{SGD favorece regiões de gradiente estável\footfullcite{smith2021origin}}
    Para SGD com \textit{batches} de tamanho $B$:
    $$\tilde{\mathcal{L}}_{SGD}(\bm{\theta}) = \mathcal{L}(\bm{\theta})
      + \frac{\eta}{4}\|\nabla \mathcal{L}\|^2
      + \frac{\eta}{4B}\sum_{b=1}^{B}\|\nabla \mathcal{L}_b - \nabla \mathcal{L}\|^2.$$
    \begin{columns}
      \begin{column}{.55\linewidth}
        \begin{itemize}
          \item O segundo termo extra corresponde à \alert{variância} do gradiente
            entre \textit{batches}.
          \item SGD favorece soluções em que o gradiente é \alert{estável} ao
            trocar de \textit{batch}.
          \item Os mínimos da função de custo original não se movem,
            mas as trajetórias e o mínimo escolhido podem mudar.
        \end{itemize}
      \end{column}
      \begin{column}{.45\linewidth}
        \centering
        \includegraphics[width=\linewidth,height=.55\textheight,keepaspectratio]{slides/img/01-mlp/reg-implicit-trajectory.png}
      \end{column}
    \end{columns}
  \end{frame}

  %======================================================================
  \section*{Referências e Leituras}
  %======================================================================

  %--- Slide: leituras ---
  \begin{frame}{Leituras Recomendadas}
    \leiturasize
    \sloppy
    \begin{itemize}
      \item Capítulo 9 de \fullcite{prince2023understanding}, livro disponível em
        \href{https://udlbook.github.io/udlbook/}{\textit{udlbook.github.io/udlbook}}.
      \item Capítulo 7 de \fullcite{goodfellow2016deep}, tratamento clássico de regularização.
      \item Seção 5.5 de \textsc{Bishop}, C.\ M. \textit{Pattern Recognition and Machine Learning}, Springer, 2006: conexão Bayesiana entre L2/L1 e \textit{priors} gaussiano/laplaciano.
      \item MNIST 1D: \citeauthoryear{greydanus2020scaling}.
      \item Double descent: \citeauthoryear{belkin2019reconciling}.
      \item Dropout: \citeauthoryear{srivastava2014dropout}; Monte Carlo
        Dropout: \citeauthoryear{gal2016dropout}.
      \item Regularização implícita: \citeauthoryear{barrett2020implicit};
        \citeauthoryear{smith2021origin}.
    \end{itemize}
  \end{frame}

  \begin{frame}[allowframebreaks]{Referências}
    \sloppy
    \printbibliography[heading=none]
  \end{frame}

\end{document}
